% Generated draft: requires mathematical and source review.
\documentclass[11pt]{article}
\usepackage[margin=25mm]{geometry}
\usepackage{fontspec}
\setmainfont{Latin Modern Roman}
\usepackage{amsmath,amssymb,mathtools,unicode-math}
\setmathfont{Latin Modern Math}
\usepackage{enumitem,needspace,xurl,hyperref}
\setlength{\parindent}{0pt}
\setlength{\parskip}{5pt}
\setlength{\emergencystretch}{3em}
\newcommand{\N}{\mathbb{N}}
\newcommand{\Z}{\mathbb{Z}}
\newcommand{\Q}{\mathbb{Q}}
\newcommand{\R}{\mathbb{R}}
\newcommand{\C}{\mathbb{C}}
\newcommand{\F}{\mathbb{F}}
\newcommand{\A}{\mathbb{A}}
\newcommand{\PP}{\mathbb{P}}
\newcommand{\E}{\mathbb{E}}
\newcommand{\eps}{\varepsilon}
\newcommand{\dd}{\,\mathrm{d}}
\begin{document}
\section*{UnsolvedMath ID: 800643}
\textbf{Source problem code:} catalogue entry
\subsection{Definitions and mathematical statement}
The supplied excerpt does not specify the ambient class of spaces or define its cover classes. The following gives a conventional open-cover interpretation, with the ambient class and the countability convention left as explicit unresolved parameters.

Fix a class \(\mathcal C\) of topological spaces and a cover convention \(d\in\{\mathrm{countable},\mathrm{arbitrary}\}\). For a topological space \(X\), let \(\mathcal O_d(X)\) consist of families \(\mathcal U\) of open subsets of \(X\) such that
\[
\bigcup\mathcal U=X,\qquad X\notin\mathcal U,
\]
with the additional requirement that \(\mathcal U\) be at most countable when \(d=\mathrm{countable}\). Define the following classes of covers within \(\mathcal O_d(X)\):
\begin{itemize}
\item \(\Omega_d(X)\) consists of the \(\omega\)-covers: for every finite set \(F\subseteq X\), some \(U\in\mathcal U\) contains \(F\).
\item \(\Gamma_d(X)\) consists of the \(\gamma\)-covers: \(\mathcal U\) is infinite and, for every \(x\in X\), the set \(\{U\in\mathcal U:x\notin U\}\) is finite.
\item \(\mathrm T_d(X)\) consists of the \(\tau\)-covers: each point of \(X\) belongs to infinitely many members of \(\mathcal U\), and for every \(x,y\in X\), at least one of the sets
\[
\{U\in\mathcal U:x\in U,\ y\notin U\},\qquad
\{U\in\mathcal U:y\in U,\ x\notin U\}
\]
is finite.
\end{itemize}
For cover classes \(\mathcal A(X)\) and \(\mathcal B(X)\), the property \(\mathsf S_1(\mathcal A,\mathcal B)\) means that for every sequence \((\mathcal U_n)_{n\geq1}\) with \(\mathcal U_n\in\mathcal A(X)\), there are choices \(U_n\in\mathcal U_n\) such that
\[
\{U_n:n\geq1\}\in\mathcal B(X).
\]
The property \(\mathsf S_{\mathrm{fin}}(\mathcal A,\mathcal B)\) means that for every such sequence there are finite subfamilies \(\mathcal F_n\subseteq\mathcal U_n\) such that
\[
\bigcup_{n\geq1}\mathcal F_n\in\mathcal B(X).
\]
The selected families in these definitions are sets of open subsets, with repetitions discarded.

\textbf{Question (supplied source, Problem 16).} For which, if any, of the seven properties
\[
\begin{gathered}
\mathsf S_1(\Omega_d,\mathrm T_d),\qquad
\mathsf S_{\mathrm{fin}}(\Omega_d,\mathrm T_d),\\
\mathsf S_1(\mathrm T_d,\Gamma_d),\qquad
\mathsf S_1(\mathrm T_d,\mathrm T_d),\qquad
\mathsf S_1(\mathrm T_d,\Omega_d),\\
\mathsf S_{\mathrm{fin}}(\mathrm T_d,\mathrm T_d),\qquad
\mathsf S_{\mathrm{fin}}(\mathrm T_d,\Omega_d)
\end{gathered}
\]
does preservation under finite powers hold? Explicitly, for each candidate property \(P\), the preservation assertion is
\[
\forall X\in\mathcal C\;\forall m\in\{1,2,3,\ldots\},\qquad
P(X)\Longrightarrow P(X^m),
\]
where \(X^m\) carries the product topology and the same cover convention is used on every power.

The intended \(\mathcal C\), the choice of \(d\), and the agreement of these conventional cover definitions with the source must be checked before this becomes an unqualified formulation of the original problem. The supplied registry labels the problem open; no independent status assessment is asserted here.
\subsection{Short English statement}
Tsaban's Problem 16 asks which, if any, of the seven listed selection properties are preserved when a space is replaced by any positive finite Cartesian power with the product topology. The supplied excerpt omits the ambient class of spaces and cover conventions, so the formulation retains these as unresolved parameters.
\subsection{Sources}
\begin{itemize}
\item Local UnsolvedMath registry, supplied complete record: permanent ID \(800643\), exact source code \texttt{catalogue entry}, title ``Selection Principles Milestone Problems --- Problem 16.'' The record attributes the problem to Tsaban, gives the proposed year as 2003, and supplies the status label ``open.'' These are supplied registry claims, not independently verified findings.
\item Tsaban, \emph{Selection Principles in Mathematics: A Milestone of Open Problems} (2003), Problem 16. Supplied source URL: \url{https://arxiv.org/abs/math/0312182}. The background records access on 2026-10-16 and extraction by \texttt{arxiv\_latex}; this is supplied provenance, not a claim of access or source verification in preparing this draft.
\end{itemize}
\end{document}
